% ==============================================================================
% Method Section — MGS-DT
% ==============================================================================

\section{Method}
\label{sec:method}

We present Multi-Granular Semantic Decision Transformer (MGS-DT), a language-enhanced offline sequential decision-making framework for ICU sepsis treatment. As illustrated in Figure~\ref{fig:architecture}, MGS-DT comprises four components: (1) a \textit{Structured-to-Language Clinical Semantic Abstraction} (S2L-CSA) module that converts structured EHR states into multi-granular clinical language descriptions; (2) a \textit{Multi-Granular Semantic Decision Transformer} ($\pi$) that predicts treatment actions by jointly modeling trajectory tokens and semantic tokens; (3) a \textit{State Predictor} ($\Omega$) that estimates post-action state transitions with calibrated uncertainty; and (4) an \textit{Uncertainty-Aware Counterfactual Self-Play} mechanism that augments offline training with high-confidence counterfactual samples via confidence-gated semi-bootstrapping.

% ------------------------------------------------------------------------------
\subsection{Problem Formulation}
\label{sec:problem}

We model the ICU sepsis treatment process as an offline Markov Decision Process (MDP) defined by the tuple $\langle \mathcal{S}, \mathcal{A}, P, R, \gamma \rangle$, where $\mathcal{S} \subset \mathbb{R}^{d_s}$ is a continuous state space of structured EHR features ($d_s = 45$), $\mathcal{A} = \{1, \ldots, |\mathcal{A}|\}$ is a discrete action space of $|\mathcal{A}| = 25$ treatment clusters, $P(s'|s,a)$ is the unknown transition dynamics, $R(s,a)$ is a reward signal derived from the SAPS-II acuity change $\Delta\text{SAPS2}$, and $\gamma \in (0,1]$ is the discount factor.

Given a dataset $\mathcal{D} = \{\tau^{(i)}\}_{i=1}^{N}$ of offline clinical trajectories, each trajectory $\tau = \{(s_t, a_t, r_t)\}_{t=0}^{T}$ records the sequence of patient states, clinician actions, and acuity outcomes observed during ICU stays. The objective is to learn a treatment policy $\pi^*$ that minimizes the cumulative SAPS-II deterioration:
\begin{equation}
    \pi^* = \arg\min_{\pi} \; \mathbb{E}_{\tau \sim \pi}\left[\sum_{t=0}^{T} \Delta\text{SAPS2}_t\right].
\end{equation}

We frame this as sequence modeling over trajectories following the Decision Transformer paradigm~\cite{chen2021decision}: the policy conditions on the return-to-go $\hat{R}_t = \sum_{t'=t}^{T} r_{t'}$, the history of states and actions, and optionally additional context, to predict the action at each timestep.

% ------------------------------------------------------------------------------
\subsection{Structured-to-Language Clinical Semantic Abstraction}
\label{sec:s2l}

A key premise of MGS-DT is that structured EHR variables, while information-rich, encode decision-relevant clinical semantics at a level too fine-grained for numerical trajectory models to exploit directly. We therefore introduce a \textit{Structured-to-Language Clinical Semantic Abstraction} (S2L-CSA) module that maps each patient state $s_t$ into three complementary language descriptions, each targeting a distinct aspect of clinical decision-making.

\paragraph{Profile semantics (who is this patient?)} The Profile view $\mathcal{V}_{\text{prof}}(s_t)$ provides a global patient snapshot oriented toward severity assessment and treatment targeting. It is generated by applying clinically-motivated threshold classifiers to the raw state vector: severity level (e.g., mild, moderate, severe, critical) derived from composite vital sign and laboratory thresholds; organ system involvement (e.g., cardiovascular, respiratory, renal, hepatic, hematologic) identified from organ-specific variable subsets; ventilation status; and a target acuity improvement expressed as cumulative SAPS-II reduction. Example output:
\begin{quote}
\textit{``CRITICAL sepsis patient requiring immediate intervention. Severity: high. Primary organ involvement: cardiovascular, respiratory. Patient is on mechanical ventilation. Target: cumulative SAPS-II improvement of 3.2 points.''}
\end{quote}

\paragraph{Context semantics (what just happened?)} The Context view $\mathcal{V}_{\text{ctx}}(s_t)$ characterizes the dynamic disease trajectory by summarizing short-term trends in vital signs (heart rate, blood pressure, respiratory rate, temperature), laboratory markers (lactate trend, creatinine, WBC, shock index, P/F ratio), fluid balance, and the most recent treatment action. It explicitly compares the current state to the preceding state to capture the direction and magnitude of clinical changes:
\begin{quote}
\textit{``After aggressive IV fluid resuscitation: blood pressure improved from hypotensive to borderline range; lactate decreased but remains elevated; urine output remains inadequate.''}
\end{quote}

\paragraph{Advisory semantics (what should we focus on?)} The Advisory view $\mathcal{V}_{\text{adv}}(s_t)$ distills the current decision context into a concise risk summary and treatment priority. It enumerates individual lab abnormalities, assigns an overall risk level, and identifies organ-specific protection priorities:
\begin{quote}
\textit{``Current risks: critically high lactate, worsening coagulation (elevated INR), declining neurological status (low GCS). Risk level: critical. Priority: hemodynamic optimization and coagulation support.''}
\end{quote}

\paragraph{Implementation.} Each view is produced by a deterministic rule-based verbalizer that applies clinical thresholds (e.g., systolic BP $< 90$ mmHg $\to$ ``hypotensive'', lactate $\geq 4.0$ mmol/L $\to$ ``critically high'') with randomized template selection (2--5 variants per category) to ensure linguistic diversity. The resulting text is encoded by a frozen pretrained language model (Qwen2.5-0.5B-Instruct) into dense embeddings:
\begin{equation}
    \mathbf{e}^{\text{prof}}_t = \text{Enc}\bigl(\mathcal{V}_{\text{prof}}(s_t)\bigr), \quad
    \mathbf{e}^{\text{ctx}}_t = \text{Enc}\bigl(\mathcal{V}_{\text{ctx}}(s_t)\bigr), \quad
    \mathbf{e}^{\text{adv}}_t = \text{Enc}\bigl(\mathcal{V}_{\text{adv}}(s_t)\bigr),
\end{equation}
where each $\mathbf{e} \in \mathbb{R}^{896}$ is the mean-pooled last hidden state. Critically, S2L-CSA operates entirely on structured EHR variables---it does \textit{not} require clinical free text or invoke the language model for generative reasoning. The language model serves purely as a fixed text encoder, converting verbalized clinical descriptions into a semantic representation space.

% ------------------------------------------------------------------------------
\subsection{Multi-Granular Semantic Decision Transformer}
\label{sec:mgsdt}

Given the semantic embeddings from S2L-CSA, MGS-DT constructs a unified token sequence that interleaves \textit{trajectory tokens} (numerical) and \textit{semantic tokens} (language-derived) within a causal Transformer architecture.

\paragraph{Token construction.} For each timestep $t$ in the context window of length $K$, MGS-DT assembles six token embeddings:
\begin{equation}
    \underbrace{\mathbf{z}^{\text{prof}}_t}_{\text{Profile}},\;
    \underbrace{\mathbf{z}^{\text{rtg}}_t}_{\text{RTG}},\;
    \underbrace{\mathbf{z}^{\text{ctx}}_t}_{\text{Context}},\;
    \underbrace{\mathbf{z}^{\text{adv}}_t}_{\text{Advisory}},\;
    \underbrace{\mathbf{z}^{s}_t}_{\text{State}},\;
    \underbrace{\mathbf{z}^{a}_t}_{\text{Action}},
\end{equation}
where each token is obtained by projecting the corresponding input into a shared embedding dimension $d$:
\begin{align}
    \mathbf{z}^{\text{prof}}_t &= W_{\text{prof}}\,\mathbf{e}^{\text{prof}}_t + b_{\text{prof}}, &
    \mathbf{z}^{\text{rtg}}_t &= W_{\text{rtg}}\,\hat{R}_t + b_{\text{rtg}}, \\
    \mathbf{z}^{\text{ctx}}_t &= W_{\text{ctx}}\,\mathbf{e}^{\text{ctx}}_t + b_{\text{ctx}}, &
    \mathbf{z}^{\text{adv}}_t &= W_{\text{adv}}\,\mathbf{e}^{\text{adv}}_t + b_{\text{adv}}, \\
    \mathbf{z}^{s}_t &= W_s\,s_t + b_s, &
    \mathbf{z}^{a}_t &= \text{Emb}(a_t),
\end{align}
followed by a $\tanh$ activation. Each projection maps from its source dimension (896 for semantic tokens, 45 for states, 1 for RTG) to the shared Transformer hidden dimension $d = 128$.

\paragraph{Interleaved sequence layout.} The full input sequence of length $6K$ tokens is arranged as:
\begin{equation}
    \bigl[\underbrace{\mathbf{z}^{\text{prof}}_0,\, \mathbf{z}^{\text{rtg}}_0,\, \mathbf{z}^{\text{ctx}}_0,\, \mathbf{z}^{\text{adv}}_0,\, \mathbf{z}^{s}_0,\, \mathbf{z}^{a}_0}_{\text{timestep 0}},\;
    \ldots,\;
    \underbrace{\mathbf{z}^{\text{prof}}_{K-1},\, \mathbf{z}^{\text{rtg}}_{K-1},\, \mathbf{z}^{\text{ctx}}_{K-1},\, \mathbf{z}^{\text{adv}}_{K-1},\, \mathbf{z}^{s}_{K-1},\, \mathbf{z}^{a}_{K-1}}_{\text{timestep } K\text{-}1}\bigr].
\end{equation}
This interleaving ensures that within each timestep, the state token $\mathbf{z}^{s}_t$ is preceded by its three semantic tokens and the return-to-go, providing rich contextual information before the action prediction position.

\paragraph{Position encoding.} We add two position signals: (1) a learnable absolute position embedding $\mathbf{p}_{\text{abs}} \in \mathbb{R}^{6K \times d}$, and (2) a timestep-level embedding $\mathbf{p}_{\text{ts}}$ that assigns the same position index to all six tokens within the same timestep, preserving the temporal structure across the interleaved layout:
\begin{equation}
    \mathbf{h}_0 = \text{Dropout}\bigl([\mathbf{z}_0, \ldots, \mathbf{z}_{6K-1}]\bigr) + \mathbf{p}_{\text{abs}} + \mathbf{p}_{\text{ts}}.
\end{equation}

\paragraph{Causal Transformer and action prediction.} The embedded sequence is processed by a stack of $L = 4$ Transformer blocks, each consisting of layer normalization, multi-head causal self-attention ($H = 8$ heads), residual connection, layer normalization, a two-layer MLP (expansion factor 4), and another residual connection. The causal mask ensures each token attends only to preceding tokens in the sequence. Action logits are extracted exclusively from the \textit{state token positions} (every 6th position starting at offset 4):
\begin{equation}
    \hat{a}_t = \text{argmax}_{a \in \mathcal{A}} \; \text{softmax}\bigl(W_{\text{out}}\,\text{LN}\bigl(\mathbf{h}^{(L)}_{6t+4}\bigr)\bigr),
\end{equation}
where $\mathbf{h}^{(L)}_{6t+4}$ denotes the hidden state at the state token position of timestep $t$ after $L$ Transformer layers. This design ensures that each action prediction conditions on all preceding trajectory and semantic information within the causal context.

% ------------------------------------------------------------------------------
\subsection{State Predictor with Calibrated Uncertainty}
\label{sec:state_predictor}

To enable counterfactual reasoning without online environment interaction, we introduce a \textit{State Predictor} $\Omega$ that estimates the outcome of a candidate action. Given the current state $s_t$ and a candidate action $a_t$, the State Predictor outputs a predicted next state mean $\boldsymbol{\mu}$, per-dimension uncertainty $\boldsymbol{\sigma}$, and predicted SAPS-II delta $\hat{\delta}$:
\begin{equation}
    (\boldsymbol{\mu},\, \boldsymbol{\sigma},\, \hat{\delta}) = \Omega(s_t,\, a_t),
\end{equation}
where $\boldsymbol{\mu}, \boldsymbol{\sigma} \in \mathbb{R}^{d_s}$ and $\hat{\delta} \in \mathbb{R}$.

\paragraph{Architecture.} The State Predictor is an MLP with three hidden layers of width 256, ReLU activations, and dropout ($p = 0.2$). The input concatenates the state vector $s_t \in \mathbb{R}^{45}$ and a one-hot action encoding $\mathbf{a}_t \in \mathbb{R}^{|\mathcal{A}|}$ (70 dimensions total). The backbone feeds three output heads:
\begin{itemize}
    \item \textbf{Mean head}: $W_{\mu} \in \mathbb{R}^{256 \times d_s}$, predicting the expected next state.
    \item \textbf{Uncertainty head}: $W_{\sigma} \in \mathbb{R}^{256 \times d_s}$, predicting $\log\boldsymbol{\sigma}$, clamped to $[-4, 4]$ for numerical stability.
    \item \textbf{Acuity head}: $W_{\delta} \in \mathbb{R}^{256 \times 1}$, predicting the SAPS-II delta.
\end{itemize}

\paragraph{Uncertainty estimation.} We employ Monte Carlo (MC) Dropout~\cite{gal2016dropout} for epistemic uncertainty quantification. At inference time, we keep dropout active and draw $M$ stochastic forward passes, estimating the total uncertainty as the sum of epistemic and aleatoric components:
\begin{equation}
    \boldsymbol{\sigma}(s_t, a_t) = \underbrace{\text{std}\bigl(\{\boldsymbol{\mu}^{(m)}\}_{m=1}^{M}\bigr)}_{\text{epistemic}} + \underbrace{\frac{1}{M}\sum_{m=1}^{M}\exp\bigl(\log\boldsymbol{\sigma}^{(m)}\bigr)}_{\text{aleatoric}},
\end{equation}
where $\boldsymbol{\mu}^{(m)}$ and $\log\boldsymbol{\sigma}^{(m)}$ are the outputs of the $m$-th MC sample. We use $M = 10$ samples in practice.

\paragraph{Training objective.} The State Predictor is trained jointly with the policy in Stage~1 using a negative log-likelihood loss for state prediction and an MSE loss for acuity prediction:
\begin{equation}
    \mathcal{L}_{\text{NLL}} = \frac{1}{d_s}\sum_{j=1}^{d_s}\left[\frac{(s_{t+1,j} - \mu_j)^2}{2\sigma_j^2 + \epsilon} + \log\sigma_j\right], \qquad
    \mathcal{L}_{\delta} = \text{MSE}(\hat{\delta},\, \Delta\text{SAPS2}_t).
\end{equation}

% ------------------------------------------------------------------------------
\subsection{Uncertainty-Aware Counterfactual Self-Play}
\label{sec:selfplay}

Offline RL suffers from distribution shift: the learned policy may propose actions not well-covered by the observed clinician trajectories. To address this while respecting offline safety constraints, we introduce a two-stage training procedure. Stage~1 performs behavioral cloning to establish a stable initialization; Stage~2 augments the training with counterfactual samples generated through uncertainty-aware self-play.

\subsubsection{Stage 1: Semantic-Augmented Behavioral Cloning}
\label{sec:stage1}

In Stage~1, both the policy $\pi$ and the State Predictor $\Omega$ are trained on the offline dataset $\mathcal{D}$. The policy minimizes a masked cross-entropy loss over clinician actions:
\begin{equation}
    \mathcal{L}_{\pi}^{\text{BC}} = -\frac{1}{\sum_t m_t}\sum_{t=0}^{T} m_t \cdot \log \pi_\theta(a_t \mid s_{<t}, a_{<t}, \hat{R}_{<t}, \mathbf{e}_{<t}),
\end{equation}
where $m_t \in \{0,1\}$ is a trajectory mask handling variable-length sequences, and $\mathbf{e}_{<t}$ denotes the history of semantic embeddings. The State Predictor is trained concurrently with $\mathcal{L}_{\text{NLL}} + \mathcal{L}_{\delta}$. Both models are optimized with AdamW using gradient clipping at 1.0.

\subsubsection{Stage 2: Counterfactual Self-Play with Confidence-Gated Semi-Bootstrap}
\label{sec:stage2}

Stage~2 iteratively searches for treatment improvements over the observed clinician behavior, using the State Predictor to evaluate counterfactual actions under uncertainty constraints.

\paragraph{Stratified action search.} At each self-play iteration, we sample a trajectory-timestep pair $(s_t, a_t^{\text{doc}}, s_{t+1})$ from $\mathcal{D}$, stratified by initial SAPS-II severity (low: $<$75, mid: 75--85, high: $\geq$85). We construct a discrete candidate set $\mathcal{C}_t = \{a \in \mathcal{A} : |a - a_t^{\text{doc}}| \leq r\} \cup \{a_t^{\pi}\}$, where $r = 2$ is the search radius and $a_t^{\pi}$ is the policy's greedy suggestion. For each candidate $a \in \mathcal{C}_t$, the State Predictor estimates the resulting state change and uncertainty.

\paragraph{Confidence filtering.} To prevent unreliable extrapolation into data-sparse regions, we reject candidates whose mean predicted uncertainty exceeds a threshold:
\begin{equation}
    \mathcal{C}_t^{\text{safe}} = \bigl\{a \in \mathcal{C}_t : \bar{\sigma}(s_t, a) < \sigma_{\text{thresh}}\bigr\},
\end{equation}
where $\bar{\sigma}$ denotes the mean uncertainty across state dimensions and $\sigma_{\text{thresh}} = 2.0$ in our experiments. Only candidates passing this filter are considered for policy improvement.

\paragraph{Advantage-guided policy update.} From the safe candidate set, we select the action with the lowest predicted SAPS-II delta:
\begin{equation}
    a_t^* = \arg\min_{a \in \mathcal{C}_t^{\text{safe}}} \hat{\delta}(s_t, a).
\end{equation}
If the counterfactual action demonstrates improvement over the clinician action, i.e., $\text{adv} = \hat{\delta}(s_t, a_t^{\text{doc}}) - \hat{\delta}(s_t, a_t^*) > 0$, the policy is updated to increase the likelihood of $a_t^*$ via MSE loss on the softmax output:
\begin{equation}
    \mathcal{L}_{\pi}^{\text{CF}} = \bigl\|\text{softmax}(\mathbf{o}_t) - \mathbf{1}_{a_t^*}\bigr\|_2^2,
\end{equation}
where $\mathbf{o}_t$ denotes the action logits at timestep $t$.

\paragraph{Confidence-gated semi-bootstrap for the State Predictor.} The State Predictor is updated using a semi-bootstrapping loss that interpolates between the real next state $s_{t+1}$ and the predictor's own counterfactual estimate $\boldsymbol{\mu}(s_t, a_t^*)$:
\begin{equation}
    \mathcal{L}_{\Omega} = \alpha_t \cdot \bigl\|\boldsymbol{\mu}(s_t, a_t^*) - \hat{\boldsymbol{\mu}}^*\bigr\|_2^2 + (1 - \alpha_t) \cdot \bigl\|\boldsymbol{\mu}(s_t, a_t^*) - s_{t+1}\bigr\|_2^2,
\end{equation}
where $\hat{\boldsymbol{\mu}}^*$ is the State Predictor's own best-candidate predicted state (self-distillation target), and the mixing coefficient $\alpha_t$ is confidence-gated:
\begin{equation}
    \alpha_t = \begin{cases}
        \alpha_{\text{base}}(t) \cdot c(s_t, a_t^*) & \text{if } c(s_t, a_t^*) \geq c_{\min}, \\
        0 & \text{otherwise},
    \end{cases}
\end{equation}
where $c(s_t, a) = 1 - \bar{\sigma}(s_t, a) / \sigma_{\text{thresh}}$ is the normalized confidence, and $\alpha_{\text{base}}(t)$ decays linearly from 1.0 to $\alpha_{\min} = 0.3$ over 50K training steps. This design ensures that the State Predictor only absorbs its own counterfactual predictions when they are sufficiently confident, preventing self-reinforcing error accumulation.

\paragraph{Total Stage~2 loss.} The combined loss for each self-play sample is:
\begin{equation}
    \mathcal{L}_{\text{total}} = \mathcal{L}_{\pi}^{\text{CF}} + \lambda_{\Omega}\,\mathcal{L}_{\Omega},
\end{equation}
where $\lambda_{\Omega} = 0.5$ balances the policy and State Predictor updates.

% ------------------------------------------------------------------------------
\subsection{Autoregressive Policy Evaluation}
\label{sec:rollout}

At evaluation time, MGS-DT generates treatment trajectories autoregressively over a horizon of $H = 10$ steps. Starting from the initial real patient state, at each step $t$:
\begin{enumerate}
    \item \textbf{Semantic generation}: Profile, Context, and Advisory descriptions are generated from the current (possibly predicted) state via S2L-CSA and encoded into embeddings.
    \item \textbf{Policy prediction}: The causal Transformer processes the growing context window (truncated to $K = 20$ timesteps) and outputs the greedy action $\hat{a}_t$.
    \item \textbf{State transition}: The State Predictor $\Omega$ estimates the next state $\boldsymbol{\mu}_{t+1}$, uncertainty $\boldsymbol{\sigma}_{t+1}$, and SAPS-II delta $\hat{\delta}_t$ via MC Dropout.
    \item \textbf{Context update}: The predicted state, action, and semantic embeddings are appended to the rolling context.
\end{enumerate}
The total treatment quality is measured as the cumulative predicted SAPS-II change $\sum_{t=0}^{H-1}\hat{\delta}_t$, with early stopping if cumulative uncertainty exceeds a safety threshold. Patients are stratified by initial SAPS-II severity for subgroup analysis.
